% © 2026 Ing. David Jaroš — CC BY-NC-ND 4.0
%
% This work is licensed under a Creative Commons Attribution-NonCommercial-NoDerivatives
% 4.0 International License (CC BY-NC-ND 4.0).
%
% B_base — Explicit a_4 Computation for ℂ⊗ℍ Spectral Triple (Approaches F1–F4)
% Track: CORE — α derivation — v64 baseline
% Date: 2026-03-08

\ifdefined\INCLUDEMODE
  % Being included in another document — skip preamble
\else
  \documentclass[11pt,a4paper]{article}
  \usepackage{amsmath,amssymb,amsthm}
  \usepackage{hyperref}
  \usepackage{booktabs}

  \newtheorem{theorem}{Theorem}
  \newtheorem{lemma}{Lemma}
  \newtheorem{proposition}{Proposition}
  \newtheorem{conjecture}{Conjecture}
  \theoremstyle{definition}
  \newtheorem{gap}{Gap}
  \theoremstyle{remark}
  \newtheorem{remark}{Remark}
  \newtheorem{observation}{Observation}

  \title{$B_{\mathrm{base}}$ — Explicit $a_4$ for the $\mathbb{C}\otimes\mathbb{H}$
    Spectral Triple\\
    \large Approaches F1--F4: UBT Triplet, Product Action,\\
    \large Wodzicki Residue, and Spectral Invariants}
  \author{Ing.\ David Jaroš}
  \date{2026-03-08 (v64)}

  \begin{document}
  \maketitle
\fi

%--------------------------------------------------------------------
\section{Recap: Where E4 Ended — and What It Missed}
\label{sec:recap_e4}
%--------------------------------------------------------------------

\noindent\textbf{No circular reasoning.}
$\alpha^{-1} = 137$ and $n^* = 137$ are \emph{never} used as inputs
anywhere in this document.
The only inputs are: $N_{\mathrm{eff}} = 12$ (proved from
$\mathbb{C}\otimes\mathbb{H}$ algebra), $\dim_\mathbb{R}(\mathrm{Im}\,\mathbb{H}) = 3$
(algebraic fact), and the UBT kinetic action $\mathcal{S}[\Theta]$.

\subsection{Established baseline (v64)}

\begin{center}
\begin{tabular}{lll}
  \toprule
  Quantity & Value & Status \\
  \midrule
  $B_0 = 2\pi N_{\mathrm{eff}}/3$ & $25.133$ & \textbf{[PROVED]} (one-loop, flat space) \\
  $B_{\mathrm{base}} = N_{\mathrm{eff}}^{3/2}$ & $41.569$ & [CONJECTURED] \\
  $B_{\mathrm{base}}/B_0 = (3/2\pi)\sqrt{N_{\mathrm{eff}}}$ & $1.6540$ & algebraic identity \\
  \bottomrule
\end{tabular}
\end{center}

\subsection{What E4 actually computed}

Approach E4 (\texttt{b\_base\_new\_directions.tex}) applied the
\emph{standard Connes--Lott} spectral triple construction for the
Standard Model:
\begin{itemize}
  \item Algebra: $A_F = \mathbb{C} \oplus \mathbb{H} \oplus \mathrm{Mat}(3,\mathbb{C})$
    (the full SM algebra, not just $\mathbb{C}\otimes\mathbb{H}$).
  \item Hilbert space: $\mathcal{H}_F = \mathbb{C}^{96}$ (all SM fermion generations,
    chiralities, and colours with $N_{\mathrm{gen}} = 3$).
  \item Dirac operator: SM Yukawa coupling matrix $D_Y$.
\end{itemize}
The $\beta$-function coefficient computed from this standard triple gave
$B_{\mathrm{SU}(2)}^{\mathrm{NCG}} = -32/3 \approx -10.67$ (negative,
asymptotically free), which is not $B_{\mathrm{base}} = +41.57$.

\subsection{The key difference E4 overlooked}

\textbf{UBT is not the Standard Model.}
The UBT field $\Theta$ takes values in $\mathbb{C}\otimes\mathbb{H} \cong
\mathrm{Mat}(2,\mathbb{C})$ — a \emph{single} biquaternionic field,
not $\mathbb{C}^{96}$ SM fermions.
The natural spectral triple for UBT is:
\begin{center}
\begin{tabular}{lll}
  \toprule
  & Standard NCG (SM, E4) & UBT (this document) \\
  \midrule
  Algebra $A$ & $\mathbb{C}\oplus\mathbb{H}\oplus\mathrm{Mat}(3,\mathbb{C})$
    & $\mathbb{C}\otimes\mathbb{H} \cong \mathrm{Mat}(2,\mathbb{C})$ \\
  Hilbert space $\mathcal{H}$ & $\mathbb{C}^{96}$
    & $\mathbb{C}^4$ (adjoint) or $\mathbb{C}^2$ (fundamental) \\
  Field & 96 SM fermions & 1 biquaternionic scalar $\Theta$ \\
  $B$ sign & $-$ (asymptotic freedom) & $+$ (IR freedom, like QED) \\
  \bottomrule
\end{tabular}
\end{center}
Approaches F1--F4 below treat the UBT spectral triple in its own right.

%--------------------------------------------------------------------
\section{Approach F1: UBT Spectral Triple — Natural Choice of
  $\mathcal{H}$ and $a_4$}
\label{sec:F1}
%--------------------------------------------------------------------

\textbf{Status: [PARTIAL] — adjoint representation is natural; exact $a_4$
computed; does not give $N_{\mathrm{eff}}^{3/2}$ directly.}

\subsection{The natural spectral triple for UBT}

Since $\Theta \in \mathbb{C}\otimes\mathbb{H}$ itself, the most natural
choice is to take $\mathcal{H}$ to be the algebra acting on itself as
a left module.
As a complex vector space,
\begin{equation}
  \mathbb{C}\otimes\mathbb{H} \;\cong\; \mathrm{Mat}(2,\mathbb{C}),
  \qquad
  \dim_\mathbb{C}(\mathbb{C}\otimes\mathbb{H}) = 4.
\end{equation}
The natural representations of $A = \mathrm{Mat}(2,\mathbb{C})$ are:
\begin{itemize}
  \item \textbf{Fundamental:} $\mathbb{C}^2$ (dim$_\mathbb{C} = 2$).
  \item \textbf{Adjoint as Lie algebra:} $\mathfrak{su}(2) \otimes_\mathbb{R}\mathbb{C}
    \cong \mathbb{C}^3$ (dim$_\mathbb{C} = 3$).
  \item \textbf{$A$ on itself (left-regular):} $\mathbb{C}^4 \cong \mathrm{Mat}(2,\mathbb{C})$
    (dim$_\mathbb{C} = 4$), decomposing as $\mathbb{C}^1_{\mathrm{triv}} \oplus
    \mathbb{C}^3_{\mathrm{adj}}$ under the $\mathrm{SU}(2)$ sub-action.
\end{itemize}
We take $\mathcal{H}_F = \mathbb{C}^4$ (the \emph{left-regular} representation
of $\mathrm{Mat}(2,\mathbb{C})$).

\subsection{Dirac operator from the UBT action}

The UBT internal Dirac operator is constructed from $\mathcal{S}[\Theta]$.
In the minimal (isotropic mass) case, the internal (finite) Dirac operator on
$\mathcal{H}_F = \mathbb{C}^4 = \mathbb{C}^2_L \oplus \mathbb{C}^2_R$ is:
\begin{equation}
  D_{\mathrm{int}} \;=\; \begin{pmatrix} 0 & M^\dagger \\ M & 0 \end{pmatrix},
  \qquad M = m\,\mathbf{1}_2 \in \mathrm{Mat}(2,\mathbb{C}),
  \label{eq:Dint}
\end{equation}
where $m$ is the internal mass scale.
Then $D_{\mathrm{int}}^2 = m^2\,\mathbf{1}_4$ (a scalar multiple of the
identity), with eigenvalues $m^2$ (multiplicity 4).

\subsection{Heat kernel coefficients for the finite part}

For a \emph{finite} spectral triple, the heat kernel is exact (not asymptotic):
\begin{equation}
  K_{\mathrm{int}}(t) \;=\; \mathrm{Tr}_{\mathcal{H}_F}\bigl(e^{-t\,D_{\mathrm{int}}^2}\bigr)
  \;=\; 4\,e^{-tm^2}
  \;=\; 4\Bigl(1 - tm^2 + \tfrac{t^2m^4}{2} - \cdots\Bigr).
  \label{eq:Kint}
\end{equation}
Reading off the ``Seeley--DeWitt-like'' coefficients of the finite triple:
\begin{align}
  a_0^{\mathrm{int}} &\;=\; \dim(\mathcal{H}_F) \;=\; 4,
    \label{eq:a0int}\\
  a_2^{\mathrm{int}} &\;=\; -\mathrm{Tr}(D_{\mathrm{int}}^2) \;=\; -4m^2,
    \label{eq:a2int}\\
  a_4^{\mathrm{int}} &\;=\; \tfrac{1}{2}\,\mathrm{Tr}(D_{\mathrm{int}}^4)
    \;=\; 2m^4.
    \label{eq:a4int}
\end{align}

\subsection{$a_4$ from the gauge connection on $\mathcal{H}_F = \mathbb{C}^4$}

For the gauge field contribution to $a_4$, the relevant term is:
\begin{equation}
  a_4 \;\supset\; \frac{1}{12}\,\mathrm{Tr}_{\mathcal{H}_F}\!\bigl(\Omega_{\mu\nu}^2\bigr)
  \;=\; \frac{1}{12}\,\mathrm{Tr}_{\mathbb{C}^4}(F_{\mu\nu}^2),
  \label{eq:a4gauge}
\end{equation}
where $F_{\mu\nu}$ acts on $\mathcal{H}_F = \mathbb{C}^4 \cong \mathbb{C}^1 \oplus
\mathbb{C}^3$ as a $\mathrm{SU}(2)$ representation.
Decomposing:
\begin{equation}
  \mathrm{Tr}_{\mathbb{C}^4}(F^2)
  \;=\; \underbrace{\mathrm{Tr}_{\mathbb{C}^1}(F^2)}_{= 0\ (\text{trivial rep})}
  + \mathrm{Tr}_{\mathbb{C}^3}(F^2).
  \label{eq:a4decomp}
\end{equation}
For the adjoint representation of $\mathrm{SU}(2)$, the Dynkin index is
$T(\mathbf{adj}) = C_2(\mathbf{adj}) = 2$, and
\begin{equation}
  \mathrm{Tr}_{\mathbf{adj}}(F^2)
  \;=\; T(\mathbf{adj})\,\mathrm{Tr}_{\mathbf{fund}}(F^2)
  \;=\; 2\,\mathrm{Tr}_{\mathbf{fund}}(F^2)
  \;=\; 2 \times 2\,\mathrm{tr}(F^2)
  \;=\; 4\,\mathrm{tr}(F^2),
  \label{eq:Tradj}
\end{equation}
where $\mathrm{tr}(F^2) \equiv \frac{1}{2}\mathrm{Tr}_{\mathbf{fund}}(F^2)$ is the
normalised trace.
Substituting into~\eqref{eq:a4gauge}:
\begin{equation}
  \bigl[a_4^{(\mathrm{gauge})}\bigr]_{\mathcal{H}_F = \mathbb{C}^4}
  \;=\; \frac{1}{12} \times 4\,\mathrm{tr}(F^2)
  \;=\; \frac{1}{3}\,\mathrm{tr}(F^2).
  \label{eq:a4gauge_final}
\end{equation}

\subsection{Comparison with standard NCG and with $N_{\mathrm{eff}}^{3/2}$}

For the standard NCG fundamental-representation choice
$\mathcal{H} = \mathbb{C}^2$:
\begin{equation}
  \bigl[a_4^{(\mathrm{gauge})}\bigr]_{\mathbb{C}^2}
  \;=\; \frac{T(\mathbf{fund})}{12}\,\mathrm{Tr}_{\mathbf{fund}}(F^2)
  \;=\; \frac{1/2}{12} \times 2\,\mathrm{tr}(F^2)
  \;=\; \frac{1}{12}\,\mathrm{tr}(F^2).
  \label{eq:a4fund}
\end{equation}
The ratio of the UBT-adjoint contribution to the standard fundamental:
\begin{equation}
  \frac{[a_4]_{\mathbb{C}^4}}{[a_4]_{\mathbb{C}^2}}
  \;=\; \frac{1/3}{1/12} \;=\; 4.
  \label{eq:ratio_F1}
\end{equation}
The $\mathcal{H}_F = \mathbb{C}^4$ internal space enhances the effective
coupling coefficient by a factor of $4$ compared to the $\mathbb{C}^2$
choice.
This implies a $\beta$-function ratio:
\begin{equation}
  \frac{B_{\mathrm{F1}}}{B_0}
  \;\stackrel{?}{=}\; 4 \times \frac{B_0}{B_0} = 4
  \quad\Rightarrow\quad
  B_{\mathrm{F1}} \stackrel{?}{=} 4 B_0 = 4 \times 25.13 = 100.5.
  \label{eq:B_F1_wrong}
\end{equation}
This is \emph{not} $N_{\mathrm{eff}}^{3/2} = 41.57$.

\begin{observation}[F1 structural result]
  \label{obs:F1}
  The left-regular representation $\mathcal{H}_F = \mathbb{C}^4$ of
  $A = \mathrm{Mat}(2,\mathbb{C})$ naturally decomposes into a trivial
  plus adjoint $\mathrm{SU}(2)$-representation, giving a $a_4$ gauge
  coefficient four times larger than the fundamental-representation choice.
  The resulting $B_{\mathrm{F1}} \approx 100$ is too large by a factor
  of $\approx 2.4$ relative to $N_{\mathrm{eff}}^{3/2} = 41.57$.
  The sign is positive (the scalar field in the adjoint contributes IR freedom,
  consistent with UBT), but the magnitude is not explained.
\end{observation}

\textbf{F1 result summary:}
\begin{itemize}
  \item $a_0^{\mathrm{int}} = 4$, $a_4^{\mathrm{int, gauge}} = \frac{1}{3}\mathrm{tr}(F^2)$.
  \item Sign: positive (IR-free), consistent with UBT.
  \item Magnitude: $B_{\mathrm{F1}} \approx 4 B_0$; does not match
    $N_{\mathrm{eff}}^{3/2}$.
  \item \textbf{[PARTIAL]:} The natural adjoint representation and sign are correct,
    but the explicit $a_4$ value does not reproduce $B_{\mathrm{base}}$.
    Next step: investigate whether a non-isotropic mass matrix $M \neq m\mathbf{1}$
    in~\eqref{eq:Dint} can reduce $[a_4]$ from $4B_0$ to $N_{\mathrm{eff}}^{3/2}$.
\end{itemize}

%--------------------------------------------------------------------
\section{Approach F2: Product Spacetime $\times$ $\mathbb{C}\otimes\mathbb{H}$
  — Total Spectral Action}
\label{sec:F2}
%--------------------------------------------------------------------

\textbf{Status: [PARTIAL] — framework correct; sign consistent;
cross-terms do not produce $N_{\mathrm{eff}}^{3/2}$.}

\subsection{Product spectral triple}

Combining 4D continuous spacetime with the finite internal space gives
the \emph{product spectral triple}:
\begin{equation}
  D_{\mathrm{total}} \;=\; D_{\mathrm{space}} \otimes \mathbf{1}_{H_F}
  + \gamma_5 \otimes D_{\mathrm{int}},
  \qquad
  \mathcal{H}_{\mathrm{total}} = L^2(M, S) \otimes \mathcal{H}_F,
  \label{eq:Dtotal}
\end{equation}
where $D_{\mathrm{space}}$ is the 4D Dirac operator on the spin bundle $S$,
and $\mathcal{H}_F = \mathbb{C}^4$ as above.

\subsection{Heat kernel factorisation}

For the product triple, the heat kernel factorises:
\begin{align}
  \mathrm{Tr}_{\mathcal{H}_{\mathrm{total}}}\!\bigl(e^{-t D_{\mathrm{total}}^2}\bigr)
  &\;=\; \mathrm{Tr}_{\mathrm{space}}\!\bigl(e^{-t D_{\mathrm{space}}^2}\bigr)
    \;\times\; K_{\mathrm{int}}(t)  \notag\\
  &\;\sim\; \Bigl[\sum_{k} t^{(k-4)/2} a_k^{\mathrm{space}}\Bigr]
    \times \Bigl[a_0^{\mathrm{int}} - t\,a_2^{\mathrm{int}} + t^2\,a_4^{\mathrm{int}} - \cdots\Bigr].
    \label{eq:product_HK}
\end{align}
Collecting the power-of-$t$ terms:
\begin{align}
  a_0^{\mathrm{total}} &\;=\; a_0^{\mathrm{space}} \times a_0^{\mathrm{int}}
    \;=\; a_0^{\mathrm{space}} \times 4, \label{eq:a0total}\\
  a_2^{\mathrm{total}} &\;=\; a_2^{\mathrm{space}} \times a_0^{\mathrm{int}}
    + a_0^{\mathrm{space}} \times (-a_2^{\mathrm{int}}),
    \label{eq:a2total}\\
  a_4^{\mathrm{total}} &\;=\; a_4^{\mathrm{space}} \times a_0^{\mathrm{int}}
    + a_2^{\mathrm{space}} \times (-a_2^{\mathrm{int}})
    + a_0^{\mathrm{space}} \times a_4^{\mathrm{int}}.
    \label{eq:a4total}
\end{align}

\subsection{Evaluation of each term on flat spacetime}

On flat 4D spacetime (Minkowski vacuum, $R_{\mu\nu\rho\sigma} = 0$):
\begin{equation}
  a_0^{\mathrm{space}} = \frac{\mathrm{vol}(M)}{(4\pi)^2},
  \qquad
  a_2^{\mathrm{space}} = 0 \quad\text{(no curvature)},
  \qquad
  a_4^{\mathrm{space}} = \frac{1}{(4\pi)^2}\frac{1}{12}\int\mathrm{tr}(F^2)\,d^4x.
  \label{eq:aspace_flat}
\end{equation}
Substituting into~\eqref{eq:a4total}:
\begin{equation}
  a_4^{\mathrm{total}} \;=\;
  \underbrace{a_4^{\mathrm{space}} \times 4}_{= 4\,a_4^{\mathrm{space}}}
  + \underbrace{0 \times (-a_2^{\mathrm{int}})}_{= 0}
  + \underbrace{a_0^{\mathrm{space}} \times 2m^4}_{\text{mass term, $\propto \Lambda^0$}}.
  \label{eq:a4total_flat}
\end{equation}
The gauge coupling coefficient in the spectral action is:
\begin{equation}
  \frac{1}{g^2_{\mathrm{eff}}} \;\propto\; f_0 \times a_4^{\mathrm{total, gauge}}
  \;=\; f_0 \times a_0^{\mathrm{int}} \times a_4^{\mathrm{space, gauge}}
  \;=\; f_0 \times 4 \times \frac{1}{12(4\pi)^2}\int\mathrm{tr}(F^2).
  \label{eq:geff_F2}
\end{equation}
The internal dimension $a_0^{\mathrm{int}} = 4$ appears as a
\emph{multiplicative factor} in the gauge coupling.

\subsection{Consequence for the $\beta$-function coefficient}

The enhancement of the gauge coupling by $a_0^{\mathrm{int}} = 4$ means
that the UBT spectral action ``sees'' the gauge fields through an effective
number of degrees of freedom equal to $4 \times N_{\mathrm{rep}}$ where
$N_{\mathrm{rep}}$ is the number of spinor components of $D_{\mathrm{space}}$.
For a 4D Weyl spinor: $N_{\mathrm{rep}} = 2$.
\begin{equation}
  N_{\mathrm{eff}}^{(\mathrm{F2})} \;=\; a_0^{\mathrm{int}} \times N_{\mathrm{rep}}
  \;=\; 4 \times 2 \;=\; 8.
  \label{eq:Neff_F2}
\end{equation}
This is less than the UBT value $N_{\mathrm{eff}} = 12$.
The cross-term $a_2^{\mathrm{space}} \times a_2^{\mathrm{int}} = 0$ on flat
spacetime.
On curved spacetime, $a_2^{\mathrm{space}} = \frac{R}{6(4\pi)^2}\int d^4x$, and
$a_2^{\mathrm{int}} = -4m^2$, giving a mass-dependent curvature coupling
\begin{equation}
  a_4^{\mathrm{total}}\big|_{\mathrm{cross}}
  \;=\; \frac{4m^2 R}{6(4\pi)^2}\int d^4x
  \quad(\text{curvature-mass cross-term}).
  \label{eq:cross_term}
\end{equation}
This is a mass-dependent correction, not a pure spectral invariant.

\begin{observation}[F2 structural result]
  \label{obs:F2}
  In the product spectral action on flat spacetime, the internal dimension
  $a_0^{\mathrm{int}} = 4$ appears as an overall multiplier of the gauge
  coupling, effectively increasing the number of gauge degrees of freedom.
  The resulting effective $N_{\mathrm{eff}}^{(\mathrm{F2})} = 8$ does not
  equal the UBT value $N_{\mathrm{eff}} = 12$.
  The cross-term between spacetime curvature and internal mass is non-zero
  on curved spacetime but vanishes in flat space and carries a free parameter $m$.
\end{observation}

\textbf{F2 result summary:}
\begin{itemize}
  \item The product structure correctly describes gauge coupling enhancement
    from internal degrees of freedom.
  \item $a_0^{\mathrm{int}} = 4$ (adjoint) enhances gauge coupling by factor 4
    vs.\ $a_0^{\mathrm{int}} = 2$ (fundamental), giving $B \approx 4B_0$.
  \item $N_{\mathrm{eff}}^{(\mathrm{F2})} = 8 \neq 12$ in this counting.
  \item \textbf{[PARTIAL]:} The product spectral action gives a sensible positive
    $B$, but neither $4B_0$ nor $2B_0$ matches $N_{\mathrm{eff}}^{3/2}$.
    The discrepancy suggests that $N_{\mathrm{eff}} = 12$ (the UBT count)
    requires input beyond the pure spectral-action product formula.
\end{itemize}

%--------------------------------------------------------------------
\section{Approach F3: Wodzicki Residue for $\mathbb{C}\otimes\mathbb{H}$}
\label{sec:F3}
%--------------------------------------------------------------------

\textbf{Status: [DEAD END] — reduces to the same $a_4$ computation.}

\subsection{Wodzicki residue definition}

For a pseudo-differential operator $P$ of order $-n$ on a compact Riemannian
manifold $M$ of dimension $n$, the \emph{Wodzicki residue} is:
\begin{equation}
  \mathrm{Res}(P)
  \;=\; \int_M \int_{|ξ|=1} \mathrm{tr}\,\sigma_{-n}(P)(x,\xi)\,d\xi\,dx,
  \label{eq:Wodzicki}
\end{equation}
where $\sigma_{-n}(P)$ is the principal symbol of order $-n$.
The Wodzicki residue is a spectral invariant — it equals (up to a constant)
the $a_{n}$ Seeley--DeWitt coefficient of the associated elliptic operator.

\subsection{Application to UBT}

For the product manifold $M^4 \times F_{\mathbb{C}\otimes\mathbb{H}}$ with
$\dim(M^4) = 4$, we take $P = D_{\mathrm{total}}^{-2}$.
The principal symbol of $D_{\mathrm{total}}^{-2}$ at order $-4$ is:
\begin{equation}
  \sigma_{-4}(D_{\mathrm{total}}^{-2})(x,\xi)
  \;=\; |\xi|^{-4} \,\mathbf{1}_{\mathcal{H}_F \otimes S},
  \label{eq:sigma4}
\end{equation}
where $S$ is the 4D spinor bundle and $\mathcal{H}_F$ is the internal space.
Integrating over the unit cosphere $|ξ| = 1$ in $\mathbb{R}^4$:
\begin{equation}
  \mathrm{Res}(D^{-2})
  \;=\; \int_M \int_{S^3} |\xi|^{-4}\,\mathrm{tr}_{\mathcal{H}_F \otimes S}
    (\mathbf{1})\,d\xi\,dx
  \;=\; \dim(\mathcal{H}_F) \times \dim(S) \times \mathrm{vol}(S^3) \times \mathrm{vol}(M).
  \label{eq:Res_result}
\end{equation}
This is proportional to $a_0^{\mathrm{int}} = \dim(\mathcal{H}_F) = 4$, the
same dimension factor that appeared in F1 and F2.
The Wodzicki residue adds \emph{no new information} beyond the heat kernel
coefficient $a_0^{\mathrm{int}}$.

\subsection{Why the Wodzicki residue cannot distinguish $N_{\mathrm{eff}}^{3/2}$}

The Wodzicki residue of $D^{-2}$ on a product space equals
$\dim(\mathcal{H}_F) \times \mathrm{vol}(M) \times (\text{universal factor})$,
which depends on $\dim(\mathcal{H}_F)$ but not on its specific algebraic
structure.
To get $N_{\mathrm{eff}}^{3/2}$ from a Wodzicki residue, one would need an
operator $P$ whose symbol carries the factor $N_{\mathrm{eff}}^{3/2}$ — which
requires knowing $N_{\mathrm{eff}}^{3/2}$ already.
This is circular.

\textbf{F3 result summary:}
\begin{itemize}
  \item $\mathrm{Res}(D^{-2}) \propto \dim(\mathcal{H}_F) = 4$ on flat space.
  \item No new invariant compared to $a_4$.
  \item \textbf{[DEAD END]:} The Wodzicki residue reduces to the same computation
    as F1/F2, and cannot produce $N_{\mathrm{eff}}^{3/2}$ without circular
    input.
\end{itemize}

%--------------------------------------------------------------------
\section{Approach F4: $N_{\mathrm{eff}}$ as a Spectral Invariant —
  Systematic Exploration}
\label{sec:F4}
%--------------------------------------------------------------------

\textbf{Status: [PARTIAL] — $N_{\mathrm{eff}}^{3/2}$ has a natural algebraic
expression in terms of $\mathbb{H}$ dimensions, but no purely spectral-algebraic
formula gives it without invoking the path-integral counting already proved.}

\subsection{Setup}

We systematically search for combinations of spectral invariants of
$A = \mathbb{C}\otimes\mathbb{H} \cong \mathrm{Mat}(2,\mathbb{C})$ that equal
$N_{\mathrm{eff}}^{3/2} = 12^{3/2} = 24\sqrt{3} \approx 41.569$.
The available invariants are:
\begin{align}
  \dim_\mathbb{R}(A) &\;=\; \dim_\mathbb{R}(\mathbb{C}\otimes\mathbb{H}) = 8,
    \label{eq:dimR_A}\\
  \dim_\mathbb{C}(A) &\;=\; \dim_\mathbb{C}(\mathrm{Mat}(2,\mathbb{C})) = 4,
    \label{eq:dimC_A}\\
  \dim_\mathbb{R}(\mathbb{H}) &\;=\; 4, \label{eq:dimR_H}\\
  \dim_\mathbb{R}(\mathrm{Im}\,\mathbb{H}) &\;=\; 3, \label{eq:dimR_ImH}\\
  N_{\mathrm{rep, fund}} &\;=\; 2 \quad\text{(fundamental of $\mathrm{SU}(2)$)},\\
  N_{\mathrm{rep, adj}} &\;=\; 3 \quad\text{(adjoint of $\mathrm{SU}(2)$)},\\
  C_2(\mathbf{adj}) &\;=\; 2, \quad T(\mathbf{fund}) = \tfrac{1}{2},
    \quad T(\mathbf{adj}) = 2 \quad\text{($\mathrm{SU}(2)$ Casimirs)}.
\end{align}

\subsection{Systematic table of candidates}

\begin{center}
\begin{tabular}{llll}
  \toprule
  Candidate combination & Explicit value & Error from $41.57$ & Verdict \\
  \midrule
  $(\dim_\mathbb{R}\,A)^{3/2} = 8^{3/2}$ & $22.63$ & $-45.6\%$ & [no] \\
  $(\dim_\mathbb{C}\,A)^{3/2} = 4^{3/2}$ & $8.00$ & $-80.8\%$ & [no] \\
  $(\dim_\mathbb{R}\,\mathbb{H})^{3/2} = 4^{3/2}$ & $8.00$ & $-80.8\%$ & [no] \\
  $\dim_\mathbb{R}\,A \times \dim_\mathbb{C}\,A = 8\times4$ & $32.00$ & $-23.0\%$ & [no] \\
  $(\dim_\mathbb{R}\,A \times \dim_\mathbb{C}\,A)^{3/4} = 32^{3/4}$ & $13.45$ & $-67.7\%$ & [no] \\
  $\dim_\mathbb{R}\,A \times \sqrt{\dim_\mathbb{C}\,A} = 8\times2$ & $16.00$ & $-61.5\%$ & [no] \\
  $(\dim_\mathbb{R}\,\mathbb{H})^2 \times \dim_\mathbb{R}(\mathrm{Im}\,\mathbb{H})^{1/2}
    = 16\sqrt{3}$ & $27.71$ & $-33.4\%$ & [no] \\
  $[\dim_\mathbb{R}\,\mathbb{H} \times (\dim_\mathbb{R}\,\mathbb{H}-1)]^{3/2}
    = (4\times3)^{3/2} = 12^{3/2}$ & $41.57$ & $0\%$ & \checkmark \\
  $(\dim_\mathbb{R}\,A)^{3/2} \times \pi / \sqrt{2} = 22.63\times\pi/\sqrt{2}$
    & $50.18$ & $+20.7\%$ & [no] \\
  $C_2(\mathbf{adj}) \times (\dim_\mathbb{R}\,A)^{3/2} = 2\times22.63$ & $45.25$ & $+8.9\%$
    & [no] \\
  \bottomrule
\end{tabular}
\end{center}

\subsection{The unique natural candidate}

The only candidate that gives exactly $N_{\mathrm{eff}}^{3/2}$ is:
\begin{equation}
  \bigl[\dim_\mathbb{R}(\mathbb{H}) \times \bigl(\dim_\mathbb{R}(\mathbb{H}) - 1\bigr)\bigr]^{3/2}
  \;=\; \bigl[4 \times 3\bigr]^{3/2}
  \;=\; \bigl[\dim_\mathbb{R}(\mathbb{H}) \times \dim_\mathbb{R}(\mathrm{Im}\,\mathbb{H})\bigr]^{3/2}
  \;=\; 12^{3/2}
  \;=\; N_{\mathrm{eff}}^{3/2}.
  \label{eq:F4_candidate}
\end{equation}
This can be rewritten using the algebraic identity
\begin{equation}
  N_{\mathrm{eff}}
  \;=\; \dim_\mathbb{R}(\mathbb{H}) \times \dim_\mathbb{R}(\mathrm{Im}\,\mathbb{H})
  \;=\; 4 \times 3,
  \label{eq:Neff_algebraic}
\end{equation}
which follows from the proved counting
$N_{\mathrm{eff}} = N_{\mathrm{phases}} \times N_{\mathrm{hel}} \times N_{\mathrm{charge}}
= 3 \times 2 \times 2 = 12$, where $N_{\mathrm{phases}} = \dim_\mathbb{R}(\mathrm{Im}\,\mathbb{H}) = 3$
is already a proved algebraic fact.

\subsection{Does this constitute a new derivation of $N_{\mathrm{eff}}^{3/2}$?}

\begin{observation}[F4 algebraic identity]
  \label{obs:F4}
  Equation~\eqref{eq:F4_candidate} is an algebraic identity, not an independent
  derivation.
  It expresses $N_{\mathrm{eff}}^{3/2}$ in terms of $\dim(\mathbb{H})$ and
  $\dim(\mathrm{Im}\,\mathbb{H})$, but both of these dimensions were already
  the inputs to the proved one-loop formula $B_0 = 2\pi N_{\mathrm{eff}}/3$.
  The exponent $3/2$ is not explained by this dimensional formula; it
  must come from the Gaussian path-integral measure on
  $\mathrm{Im}(\mathbb{H})$ (Approach A2 in \texttt{b\_base\_hausdorff.tex}).
\end{observation}

A purely spectral/algebraic formula that derives the exponent $3/2$ from
the structure of $\mathbb{C}\otimes\mathbb{H}$ (without invoking the
Gaussian measure) does not emerge from this systematic search.

\textbf{F4 result summary:}
\begin{itemize}
  \item The unique natural candidate is $[\dim_\mathbb{R}(\mathbb{H}) \times
    \dim_\mathbb{R}(\mathrm{Im}\,\mathbb{H})]^{3/2} = N_{\mathrm{eff}}^{3/2}$.
  \item This is an algebraic identity expressing what is already known.
  \item \textbf{[PARTIAL]:} The algebraic structure uniquely selects
    $N_{\mathrm{eff}}^{3/2}$ among dimension combinations, but the $3/2$
    exponent itself requires the Gaussian path-integral argument (Gap~(a) /
    Gap~(b) in \texttt{b\_base\_hausdorff.tex}).
\end{itemize}

%--------------------------------------------------------------------
\section{Summary: All Approaches F1--F4}
\label{sec:summary_F}
%--------------------------------------------------------------------

\begin{center}
\begin{tabular}{lll}
  \toprule
  Approach & Result & Status \\
  \midrule
  F1: UBT spectral triple, $\mathcal{H} = \mathbb{C}^4$ (adjoint)
    & $a_4 = \frac{1}{3}\mathrm{tr}(F^2)$; $B_{\mathrm{F1}} \approx 4B_0$
    & \textbf{[PARTIAL]} \\
  F2: Product spacetime $\times$ $\mathbb{C}\otimes\mathbb{H}$
    & $a_0^{\mathrm{int}} = 4$ multiplies $B$; $N_{\mathrm{eff}}^{(\mathrm{F2})} = 8$
    & \textbf{[PARTIAL]} \\
  F3: Wodzicki residue
    & Reduces to $a_4$; no new invariant
    & \textbf{[DEAD END]} \\
  F4: $N_{\mathrm{eff}}$ as spectral invariant
    & Unique: $[\dim_\mathbb{R}(\mathbb{H})\times\dim_\mathbb{R}(\mathrm{Im}\,\mathbb{H})]^{3/2}$;
      algebraic identity only
    & \textbf{[PARTIAL]} \\
  \bottomrule
\end{tabular}
\end{center}

\subsection{Cumulative status of all approaches (A1--F4)}

\begin{center}
\begin{tabular}{llll}
  \toprule
  Approach & Mechanism & Status & File \\
  \midrule
  A1 & KK mode sum & [DEAD END] & \texttt{b\_base\_spinor\_approach.tex} \\
  A2 & Spectral zeta / Hausdorff & [MOTIVATED CONJECTURE] & \texttt{b\_base\_hausdorff.tex} \\
  A3 & Gauge orbit volume & [DEAD END] & \texttt{b\_base\_hausdorff.tex} \\
  A4 & Effective dimension & [STRUCTURAL INSIGHT] & tools \\
  C1 & Seeley--DeWitt (pert.) & [DEAD END] & \texttt{b\_base\_delta\_d.tex} \\
  C2 & Mode pair coherence & [DEAD END] & \texttt{b\_base\_delta\_d.tex} \\
  D1 & Unitarity constraint & [DEAD END] & \texttt{b\_base\_nonpert.tex} \\
  D2 & Dimensional transmutation & [DEAD END] & \texttt{b\_base\_nonpert.tex} \\
  D3 & Cartan--Killing metric & [DEAD END] & \texttt{b\_base\_nonpert.tex} \\
  E1 & Instantons on $\mathrm{SU}(2)$ & [DEAD END] & \texttt{b\_base\_new\_directions.tex} \\
  E2 & Index theorem / anomaly & [DEAD END] & \texttt{b\_base\_new\_directions.tex} \\
  E3 & Holomorphic factorisation & [DEAD END/CONFIRM] & \texttt{b\_base\_new\_directions.tex} \\
  E4 & NCG spectral triple (SM) & [PARTIAL] & \texttt{b\_base\_new\_directions.tex} \\
  F1 & UBT spectral triple, adjoint & [PARTIAL] & this file \\
  F2 & Product spacetime action & [PARTIAL] & this file \\
  F3 & Wodzicki residue & [DEAD END] & this file \\
  F4 & Spectral invariant search & [PARTIAL] & this file \\
  \bottomrule
\end{tabular}
\end{center}

\subsection{Remaining gaps and next steps}

\begin{gap}[Gap (a): Explicit $\det(\mathcal{S}'')$]
  \label{gap:a}
  The functional determinant of the second variation of the UBT action
  $\mathcal{S}[\Theta]$ on $\mathrm{Im}(\mathbb{H})$ has not been computed
  explicitly.
  F1--F4 confirm that the natural spectral dimension is $a_0^{\mathrm{int}} = 4$
  (adjoint representation), not directly $N_{\mathrm{eff}}^{3/2}$.
  Status: \textbf{[OPEN]}.
\end{gap}

\begin{gap}[Gap (b): Protection of the $3/2$ exponent]
  \label{gap:b}
  The exponent $3/2$ in $B_{\mathrm{base}} = N_{\mathrm{eff}}^{3/2}$ is
  structurally forced by the real Gaussian integration over the odd-dimensional
  space $\mathrm{Im}(\mathbb{H})$, structurally confirmed in E3
  (holomorphic factorisation, \texttt{b\_base\_new\_directions.tex} §E3)
  and motivated by A2 (Hausdorff / Gaussian path integral,
  \texttt{b\_base\_hausdorff.tex} §2).
  Approaches F1--F4 are consistent with this but do not add new protection.
  Status: \textbf{[PARTIAL PROGRESS from E3; no new progress in F1--F4]}.
\end{gap}

\begin{gap}[Gap (F): Non-isotropic Dirac operator]
  \label{gap:F}
  F1 showed that the isotropic Dirac operator $D_{\mathrm{int}} = m\mathbf{1}$
  gives $B_{\mathrm{F1}} = 4B_0 \neq N_{\mathrm{eff}}^{3/2}$.
  The next natural step is to compute $a_4$ for a \emph{non-isotropic}
  $D_{\mathrm{int}}$ (with $\mathrm{SU}(2)$-breaking mass terms), and check
  whether a specific mass matrix that respects the $\mathrm{Im}(\mathbb{H})$
  structure gives the correct value.
  Status: \textbf{[OPEN — concrete next direction]}.
\end{gap}

%--------------------------------------------------------------------
\section{In Plain Language}
\label{sec:plain}
%--------------------------------------------------------------------

\textbf{In plain language:}
We tried four new approaches to explain why the UBT running constant
$B_{\mathrm{base}} = 41.57$ instead of the proved one-loop value $B_0 = 25.13$.
All four approaches use the idea that the UBT field $\Theta$ is a biquaternionic
object ($2 \times 2$ complex matrix), and we asked what the natural quantum
calculation looks like when we treat this properly.

\textbf{F1 (UBT spectral triple):}
The most natural mathematical setting for the biquaternion field gives a
``quantum stage'' (Hilbert space) of 4 complex dimensions rather than the 2 used
by the Standard Model approach in E4.
We computed the relevant quantity ($a_4$, the Seeley--DeWitt coefficient)
exactly for this 4-dimensional stage.
The result is that the coupling constant is multiplied by 4 compared to the
minimal case, giving $B_{\mathrm{F1}} \approx 100$ — too large by a factor of 2.4.
The right sign (the theory is IR-free like electromagnetism, not asymptotically
free like QCD) is confirmed, which is an important consistency check.
But the exact value $41.57$ is not reproduced.

\textbf{F2 (Product action):}
We combined the 4D spacetime calculation with the biquaternionic internal
structure using a product formula.
The internal space contributes an overall factor (its dimension, $= 4$) to the
gauge coupling.
On flat spacetime the mixing terms (between spacetime curvature and the mass
of the field) vanish, so there are no surprises.
The effective number of modes this formula counts is 8, not the required 12,
so $B_{\mathrm{base}} = 41.57$ is not explained.

\textbf{F3 (Wodzicki residue):}
There is a sophisticated mathematical tool called the Wodzicki residue that
is supposed to be a ``topological invariant'' — a number that does not depend
on any free parameters.
We asked whether this tool gives $41.57$ directly.
The answer is no: the Wodzicki residue of the natural operator reduces to
exactly the same calculation as in F1 and F2 (it counts the dimension of the
quantum stage, which is 4).
No new information is obtained.

\textbf{F4 (Spectral invariant search):}
We systematically listed all natural combinations of dimensions and Casimir
invariants of the biquaternion algebra $\mathbb{C}\otimes\mathbb{H}$ and
checked which ones equal $41.57$.
Only one combination works:
$[\dim(\mathbb{H})\times\dim(\mathrm{Im}\,\mathbb{H})]^{3/2} = (4\times 3)^{3/2} = 41.57$.
This is the unique purely algebraic expression.
However, it is just a restatement of the already-known fact that
$N_{\mathrm{eff}} = 12$ combined with the Gaussian exponent $3/2$.
It does not explain \emph{where the $3/2$ comes from} — that question still
requires the Gaussian path-integral argument established in earlier work
(Approach A2, \texttt{b\_base\_hausdorff.tex} §2, and structural
confirmation E3, \texttt{b\_base\_new\_directions.tex} §E3).

\textbf{Overall summary of F1--F4:}
F3 is a dead end (no new information).
F1, F2, and F4 are all partial: they confirm the correct sign and algebraic
structure, but do not independently derive $B_{\mathrm{base}} = 41.57$.
The most concrete next step (Gap F) is to compute $a_4$ for a non-isotropic
Dirac operator that respects the $\mathrm{Im}(\mathbb{H})$ structure, and
check whether this shifts $B$ from $4B_0$ to $N_{\mathrm{eff}}^{3/2}$.
The fundamental obstacle remains: the $3/2$ exponent requires the Gaussian
path-integral argument, and no purely algebraic combination of spectral
invariants of $\mathbb{C}\otimes\mathbb{H}$ can generate it from scratch.

\ifdefined\INCLUDEMODE\else
\end{document}
\fi
